% ============================================================
%  EXAMEN FINAL — CURSO 4TO AÑO DE EDUCACIÓN MEDIA GENERAL
%  Duración: 90 minutos — Puntaje total: 100 puntos
%  Temas: Matemática (Ángulos y Razones Trigonométricas),
%         Física (Unidades, Vectores y Cinemática),
%         Química (Estructura Atómica y Estequiometría)
%  Autor: Gabriel Astudillo
%  Clase: 18 (sesión de evaluación)
% ============================================================

\documentclass[12pt, a4paper]{article}

% ── Paquetes ─────────────────────────────────────────────────

\usepackage{mathwizards-guia}
\mwguia{Examen Final — 4to Año}{Gabriel Astudillo $\cdot$ 90 minutos}

% ── Comandos propios de este examen ─────────────────
\providecommand{\sen}{\sin}

\begin{document}
% ============================================================

% ── Portada / Título ─────────────────────────────────────────
\mwportada{Examen Final}{Curso 4to Año de Educación Media General}{Matemática, Física y Química}

\vspace{4pt}
\begin{cuadroinfo}
\small
\textbf{Nombre:} \rule{0.55\linewidth}{0.4pt} \quad \textbf{Fecha:} \rule{0.15\linewidth}{0.4pt} \\[6pt]
\textbf{Tiempo disponible:} 90 minutos \quad$\bullet$\quad \textbf{Puntaje total:} 100 puntos \\[4pt]
\textbf{Instrucciones:}
\begin{enumerate}[leftmargin=*]
  \item Lee cada ejercicio con calma antes de resolverlo.
  \item Muestra \emph{todos} tus pasos: se evalúa cómo llegas a la respuesta, no solo la respuesta.
  \item En física y química, \textbf{no olvides las unidades} en tus resultados.
  \item No se permite el uso de calculadora.
  \item Escribe con letra clara y revisa tu examen antes de entregarlo.
\end{enumerate}
\end{cuadroinfo}

\vspace{8pt}

% ============================================================
\section{Matemática: Ángulos y Razones Trigonométricas (35 puntos)}
% ============================================================

\begin{enumerate}[leftmargin=*]
  \item \textbf{(5 pts)} Clasifica cada ángulo según su medida y calcula su complemento o suplemento:
  \begin{enumerate}[label=\alph*)]
    \item $35^{\circ}$ (clasifícalo y halla su complemento).
    \item $145^{\circ}$ (clasifícalo y halla su suplemento).
  \end{enumerate}

  \item \textbf{(5 pts)} Conversiones y cuadrantes:
  \begin{enumerate}[label=\alph*)]
    \item Convierte $60^{\circ}$ a radianes.
    \item Convierte $\dfrac{\pi}{3}$ rad a grados.
    \item Determina el cuadrante de $240^{\circ}$.
  \end{enumerate}

  \item \textbf{(5 pts)} En el triángulo rectángulo con catetos $3$ y $4$ (hipotenusa $5$), escribe el seno, el coseno y la tangente del ángulo agudo $\alpha$ cuyo cateto opuesto es $3$.

  \item \textbf{(5 pts)} Ángulos notables:
  \begin{enumerate}[label=\alph*)]
    \item La hipotenusa de un triángulo rectángulo mide $12$ y uno de sus ángulos agudos mide $30^{\circ}$. ¿Cuánto mide el cateto opuesto a ese ángulo?
    \item Calcula: $\tan 60^{\circ} \cdot \cos 30^{\circ}$.
  \end{enumerate}

  \item \textbf{(5 pts)} Desde un punto a $50$ m de la base de un edificio, el ángulo de elevación hacia la azotea es de $60^{\circ}$. ¿Qué altura tiene el edificio? (Usa $\tan 60^{\circ} = \sqrt{3}$.)

  \item \textbf{(10 pts)} Una escalera de $10$ m apoyada contra una pared forma un ángulo de $60^{\circ}$ con el suelo.
  \begin{enumerate}[label=\alph*)]
    \item ¿A qué altura de la pared llega la escalera? (Usa $\sen 60^{\circ} = \dfrac{\sqrt{3}}{2}$.)
    \item ¿A qué distancia de la pared está la base de la escalera? (Usa $\cos 60^{\circ} = \dfrac{1}{2}$.)
  \end{enumerate}
\end{enumerate}

% ============================================================
\section{Física: Unidades, Vectores y Cinemática (35 puntos)}
% ============================================================

\begin{enumerate}[leftmargin=*]
  \setcounter{enumi}{6}
  \item \textbf{(5 pts)} Clasifica en magnitud fundamental o derivada e indica su unidad del SI: $a)$ masa \quad $b)$ velocidad \quad $c)$ tiempo \quad $d)$ densidad.

  \item \textbf{(5 pts)} Escribe en notación científica: $a)\ 4\,500\,000 \qquad b)\ 0{,}00025$.

  \item \textbf{(5 pts)} Transformación de unidades:
  \begin{enumerate}[label=\alph*)]
    \item $250$ cm a metros.
    \item $72$ km/h a m/s.
    \item $15$ m/s a km/h.
  \end{enumerate}

  \item \textbf{(5 pts)} Un vector $\vec{v}$ tiene módulo $10$ y forma un ángulo de $60^{\circ}$ con el eje $X$ positivo. Halla sus componentes rectangulares. (Usa $\cos 60^{\circ} = \tfrac{1}{2}$ y $\sen 60^{\circ} = \tfrac{\sqrt{3}}{2}$.)

  \item \textbf{(5 pts)} Una persona camina $3$ km hacia el este y luego $4$ km hacia el norte. ¿Cuál es el módulo de su desplazamiento total? ¿Es lo mismo que la distancia recorrida? Explica brevemente.

  \item \textbf{(5 pts)} Un auto pasa de $10$ m/s a $30$ m/s en $5$ s:
  \begin{enumerate}[label=\alph*)]
    \item ¿Cuál es su aceleración media?
    \item ¿Qué distancia recorre en esos $5$ segundos? (Usa $\Delta x = v_i t + \tfrac{1}{2} a t^2$.)
  \end{enumerate}

  \item \textbf{(5 pts)} Una piedra se deja caer desde un puente y tarda $3$ s en llegar al agua ($g = 9{,}8$ m/s$^2$):
  \begin{enumerate}[label=\alph*)]
    \item ¿Qué altura tiene el puente? (Usa $h = \tfrac{1}{2} g t^2$.)
    \item ¿Con qué velocidad llega al agua? (Usa $v = g t$.)
  \end{enumerate}
\end{enumerate}

% ============================================================
\section{Química: Estructura Atómica y Estequiometría\\(30 puntos)}
% ============================================================

\begin{enumerate}[leftmargin=*]
  \setcounter{enumi}{13}
  \item \textbf{(5 pts)} El sodio tiene $Z = 11$ y $A = 23$:
  \begin{enumerate}[label=\alph*)]
    \item ¿Cuántos protones, neutrones y electrones tiene el átomo neutro?
    \item Escribe su notación ${}^{A}_{Z}\text{X}$.
    \item ¿Cuántos electrones tiene el ion $\text{Na}^{+}$?
  \end{enumerate}

  \item \textbf{(5 pts)} El cloro tiene $Z = 17$:
  \begin{enumerate}[label=\alph*)]
    \item Escribe su configuración electrónica.
    \item ¿Cuántos electrones de valencia tiene?
    \item ¿En qué periodo y grupo de la tabla periódica está?
  \end{enumerate}

  \item \textbf{(5 pts)} Determina los números cuánticos $(n, l, m_l, m_s)$ del último electrón del fósforo ($Z = 15$), aplicando la regla de Hund.

  \item \textbf{(5 pts)} El mol y la masa molar ($\text{Na} = 23$, $\text{H} = 1$, $\text{S} = 32$, $\text{O} = 16$):
  \begin{enumerate}[label=\alph*)]
    \item ¿Cuántos moles hay en $46$ g de sodio?
    \item Calcula la masa molar del ácido sulfúrico, $\text{H}_2\text{SO}_4$.
  \end{enumerate}

  \item \textbf{(5 pts)} Balancea las siguientes ecuaciones químicas:
  \begin{enumerate}[label=\alph*)]
    \item $\text{N}_2 + \text{H}_2 \rightarrow \text{NH}_3$
    \item $\text{CH}_4 + \text{O}_2 \rightarrow \text{CO}_2 + \text{H}_2\text{O}$
  \end{enumerate}

  \item \textbf{(5 pts)} Composición porcentual y cálculos ($\text{C} = 12$, $\text{O} = 16$, $\text{H} = 1$):
  \begin{enumerate}[label=\alph*)]
    \item Calcula el porcentaje de carbono y de oxígeno en el $\text{CO}_2$ ($M = 44$ g/mol).
    \item En la reacción $2\,\text{H}_2 + \text{O}_2 \rightarrow 2\,\text{H}_2\text{O}$: ¿cuántos gramos de agua se producen a partir de $4$ g de hidrógeno? ($M(\text{H}_2\text{O}) = 18$ g/mol.)
  \end{enumerate}
\end{enumerate}

\vspace{14pt}
\begin{center}
  {\Large\bfseries\color{mwprimario} ¿Mucho éxito!}\\[4pt]
  {\large\color{gray} Respira profundo, lee con calma y confía en todo lo que practicaste en las seis guías.}
\end{center}

\end{document}
